\documentclass{article}
\usepackage[margin=1in]{geometry}
\usepackage{booktabs}
\usepackage{hyperref}
\title{PelicanBench: A Longitudinal, Compositional, and Agentic Visual Evaluation Suite}
\author{Dylan Mordaunt}
\date{}
\begin{document}
\maketitle
\begin{abstract}
PelicanBench converts the recurring challenge of generating a pelican riding a bicycle into a dynamic evaluation suite for vector generation, compositional relations, repair, and agentic drawing. We describe typed animal, mobility-system, and interface ontologies; a rights-aware longitudinal archive; secure SVG analysis; multidimensional scorecards; human calibration; factorial contamination analysis; and trajectory evaluation. This repository release is an MVP methods and software artifact rather than a completed empirical leaderboard.
\end{abstract}
\section{Introduction}
The public anchor is memorable but exposed. Our primary benchmark samples related animals, mobile non-living objects, roles, views, and styles while preserving longitudinal comparability.
\section{Methods}
\subsection{Task ontology} Separate animal and object ontologies are composed through a typed interface ontology.
\subsection{Scoring} Source integrity, anatomy, mechanics, interaction, composition, vector quality, and instruction coverage are reported separately with critical gates.
\subsection{Agentic trajectories} We analyse improvement, area under the score trajectory, regressions, recovery, repeated actions, cost, and edit preservation.
\subsection{Human calibration} Pairwise blinded ratings calibrate automatic structural and semantic scorers.
\section{Planned analyses}
Longitudinal trends, dynamic challenge performance, repair success, application transfer, and the Pelicanmaxxing interaction index are preregistered before full evaluation.
\section{Limitations}
Historical data are observational and rights-constrained; V1 structural scores are not complete semantic judgements; empirical human calibration remains future work.
\end{document}
